\documentclass{article}
\usepackage[margin=1in]{geometry}
\usepackage{hyperref}
\usepackage{xcolor}
\usepackage{listings}
\usepackage{longtable}
\usepackage{booktabs}
\usepackage{enumitem}
\usepackage{titlesec}
\usepackage{fancyhdr}

\hypersetup{
    colorlinks=true,
    linkcolor=blue,
    citecolor=blue,
    urlcolor=blue
}

\pagestyle{fancy}
\fancyhf{}
\fancyhead[L]{RFC 8827}
\fancyhead[R]{WebRTC Security Architecture}
\fancyfoot[C]{\thepage}

\title{WebRTC Security Architecture}
\author{E. Rescorla \\ Mozilla}
\date{January 2021}

\begin{document}

\maketitle

\begin{abstract}
This document defines the security architecture for WebRTC, a
protocol suite intended for use with real-time applications that can
be deployed in browsers -- ``real-time communication on the Web''.
\end{abstract}

\noindent\textbf{Status of This Memo}

This is an Internet Standards Track document.

This document is a product of the Internet Engineering Task Force
(IETF).  It represents the consensus of the IETF community.  It has
received public review and has been approved for publication by the
Internet Engineering Steering Group (IESG).  Further information on
Internet Standards is available in Section 2 of RFC 7841.

Information about the current status of this document, any errata,
and how to provide feedback on it may be obtained at
\url{https://www.rfc-editor.org/info/rfc8827}.

\bigskip
\noindent\textbf{Copyright Notice}

Copyright (c) 2021 IETF Trust and the persons identified as the
document authors.  All rights reserved.

This document is subject to BCP 78 and the IETF Trust's Legal
Provisions Relating to IETF Documents
(\url{https://trustee.ietf.org/license-info}) in effect on the date of
publication of this document.  Please review these documents
carefully, as they describe your rights and restrictions with respect
to this document.  Code Components extracted from this document must
include Simplified BSD License text as described in Section 4.e of
the Trust Legal Provisions and are provided without warranty as
described in the Simplified BSD License.

This document may contain material from IETF Documents or IETF
Contributions published or made publicly available before November
10, 2008.  The person(s) controlling the copyright in some of this
material may not have granted the IETF Trust the right to allow
modifications of such material outside the IETF Standards Process.
Without obtaining an adequate license from the person(s) controlling
the copyright in such materials, this document may not be modified
outside the IETF Standards Process, and derivative works of it may
not be created outside the IETF Standards Process, except to format
it for publication as an RFC or to translate it into languages other
than English.

\tableofcontents

\section{Introduction}

The Real-Time Communications on the Web (RTCWEB) Working Group
standardized protocols for real-time communications between Web
browsers, generally called ``WebRTC'' [RFC8825].  The major use cases
for WebRTC technology are real-time audio and/or video calls, Web
conferencing, and direct data transfer.  Unlike most conventional
real-time systems (e.g., SIP-based [RFC3261] soft phones), WebRTC
communications are directly controlled by some Web server, via a
JavaScript (JS) API as shown in Figure~\ref{fig:simple-webrtc}.

\begin{figure}[ht]
\begin{verbatim}
                       +----------------+
                       |                |
                       |   Web Server   |
                       |                |
                       +----------------+
                           ^        ^
                          /          \
                  HTTP   /            \   HTTP
                        /              \
                       /                \
                      v                  v
                   JS API              JS API
             +-----------+            +-----------+
             |           |    Media   |           |
             |  Browser  |<---------->|  Browser  |
             |           |            |           |
             +-----------+            +-----------+
\end{verbatim}
\caption{A Simple WebRTC System}
\label{fig:simple-webrtc}
\end{figure}

A more complicated system might allow for inter-domain calling, as
shown in Figure~\ref{fig:multidomain-webrtc}.  The protocol to be used between the domains is
not standardized by WebRTC, but given the installed base and the form
of the WebRTC API is likely to be something SDP-based like SIP or
something like the Extensible Messaging and Presence Protocol (XMPP)
[RFC6120].

\begin{figure}[ht]
\begin{verbatim}
    +--------------+                +--------------+
    |              | SIP, XMPP, ... |              |
    |  Web Server  |<-------------->|  Web Server  |
    |              |                |              |
    +--------------+                +--------------+
           ^                                ^
           |                                |
     HTTP  |                                |  HTTP
           |                                |
           v                                v
           JS API                       JS API
     +-----------+                     +-----------+
     |           |        Media        |           |
     |  Browser  |<------------------->|  Browser  |
     |           |                     |           |
     +-----------+                     +-----------+
\end{verbatim}
\caption{A Multidomain WebRTC System}
\label{fig:multidomain-webrtc}
\end{figure}

This system presents a number of new security challenges, which are
analyzed in [RFC8826].  This document describes a security
architecture for WebRTC which addresses the threats and requirements
described in that document.

\section{Terminology}

The key words ``MUST'', ``MUST NOT'', ``REQUIRED'', ``SHALL'', ``SHALL NOT'',
``SHOULD'', ``SHOULD NOT'', ``RECOMMENDED'', ``NOT RECOMMENDED'', ``MAY'', and
``OPTIONAL'' in this document are to be interpreted as described in
BCP~14 [RFC2119] [RFC8174] when, and only when, they appear in all
capitals, as shown here.

\section{Trust Model}

The basic assumption of this architecture is that network resources
exist in a hierarchy of trust, rooted in the browser, which serves as
the user's Trusted Computing Base (TCB).  Any security property which
the user wishes to have enforced must be ultimately guaranteed by the
browser (or transitively by some property the browser verifies).
Conversely, if the browser is compromised, then no security
guarantees are possible.  Note that there are cases (e.g., Internet
kiosks) where the user can't really trust the browser that much.  In
these cases, the level of security provided is limited by how much
they trust the browser.

Optimally, we would not rely on trust in any entities other than the
browser.  However, this is unfortunately not possible if we wish to
have a functional system.  Other network elements fall into two
categories: those which can be authenticated by the browser and thus
can be granted permissions to access sensitive resources, and those
which cannot be authenticated and thus are untrusted.

\subsection{Authenticated Entities}

There are two major classes of authenticated entities in the system:

\begin{description}[style=nextline]
\item[Calling services:] Web sites whose origin we can verify (optimally
  via HTTPS, but in some cases because we are on a topologically
  restricted network, such as behind a firewall, and can infer
  authentication from firewall behavior).

\item[Other users:] WebRTC peers whose origin we can verify
  cryptographically (optimally via DTLS-SRTP).
\end{description}

Note that merely being authenticated does not make these entities
trusted.  For instance, just because we can verify that
\url{https://www.example.org/} is owned by Dr.\ Evil does not mean that we
can trust Dr.\ Evil to access our camera and microphone.  However, it
gives the user an opportunity to determine whether they wish to trust
Dr.\ Evil or not; after all, if they desire to contact Dr.\ Evil
(perhaps to arrange for ransom payment), it's safe to temporarily
give them access to the camera and microphone for the purpose of the
call, but they don't want Dr.\ Evil to be able to access their camera
and microphone other than during the call.  The point here is that we
must first identify other elements before we can determine whether
and how much to trust them.  Additionally, sometimes we need to
identify the communicating peer before we know what policies to
apply.

\subsection{Unauthenticated Entities}

Other than the above entities, we are not generally able to identify
other network elements; thus, we cannot trust them.  This does not
mean that it is not possible to have any interaction with them, but
it means that we must assume that they will behave maliciously and
design a system which is secure even if they do so.

\section{Overview}

This section describes a typical WebRTC session and shows how the
various security elements interact and what guarantees are provided
to the user.  The example in this section is a ``best case'' scenario
in which we provide the maximal amount of user authentication and
media privacy with the minimal level of trust in the calling service.
Simpler versions with lower levels of security are also possible and
are noted in the text where applicable.  It's also important to
recognize the tension between security (or performance) and privacy.
The example shown here is aimed towards settings where we are more
concerned about secure calling than about privacy, but as we shall
see, there are settings where one might wish to make different
tradeoffs -- this architecture is still compatible with those
settings.

For the purposes of this example, we assume the topology shown in the
figures below.  This topology is derived from the topology shown in
Figure~\ref{fig:simple-webrtc}, but separates Alice's and Bob's identities from the process
of signaling.  Specifically, Alice and Bob have relationships with
some Identity Provider (IdP) that supports a protocol (such as OpenID
Connect) that can be used to demonstrate their identity to other
parties.  For instance, Alice might have an account with a social
network which she can then use to authenticate to other Web sites
without explicitly having an account with those sites; this is a
fairly conventional pattern on the Web. Section~\ref{sec:trust-relationships} provides an
overview of IdPs and the relevant terminology.  Alice and Bob might
have relationships with different IdPs as well.  Note: The IdP
mechanism described here has not seen wide adoption.  See Section~\ref{sec:web-based-peer-auth}
for more on the status of IdP-based authentication.

This separation of identity provision and signaling isn't
particularly important in ``closed world'' cases where Alice and Bob
are users on the same social network and have identities based on
that domain (Figure~\ref{fig:call-with-idp}).  However, there are important settings where
that is not the case, such as federation (calls from one domain to
another; see Figure~\ref{fig:federated-call-with-idp}) and calling on untrusted sites, such as where
two users who have a relationship via a given social network want to
call each other on another, untrusted, site, such as a poker site.

Note that the servers themselves are also authenticated by an
external identity service, the SSL/TLS certificate infrastructure
(not shown).  As is conventional in the Web, all identities are
ultimately rooted in that system.  For instance, when an IdP makes an
identity assertion, the Relying Party consuming that assertion is
able to verify because it is able to connect to the IdP via HTTPS.

\begin{figure}[ht]
\begin{verbatim}
                          +----------------+
                          |                |
                          |     Signaling  |
                          |     Server     |
                          |                |
                          +----------------+
                              ^        ^
                             /          \
                     HTTPS  /            \   HTTPS
                           /              \
                          /                \
                         v                  v
                      JS API              JS API
                +-----------+            +-----------+
                |           |    Media   |           |
          Alice |  Browser  |<---------->|  Browser  | Bob
                |           | (DTLS+SRTP)|           |
                +-----------+            +-----------+
                      ^      ^--+     +--^     ^
                      |         |     |        |
                      v         |     |        v
                +-----------+   |     |  +-----------+
                |           |<--------+  |           |
                |   IdP1    |   |        |    IdP2   |
                |           |   +------->|           |
                +-----------+            +-----------+
\end{verbatim}
\caption{A Call with IdP-Based Identity}
\label{fig:call-with-idp}
\end{figure}

Figure~\ref{fig:federated-call-with-idp} shows essentially the same calling scenario but with a call
between two separate domains (i.e., a federated case), as in
Figure~\ref{fig:multidomain-webrtc}.  As mentioned above, the domains communicate by some
unspecified protocol, and providing separate signaling and identity
allows for calls to be authenticated regardless of the details of the
inter-domain protocol.

\begin{figure}[ht]
\begin{verbatim}
  +----------------+    Unspecified    +----------------+
  |                |      protocol     |                |
  |    Signaling   |<----------------->|    Signaling   |
  |    Server      |  (SIP, XMPP, ...) |    Server      |
  |                |                   |                |
  +----------------+                   +----------------+
           ^                                   ^
           |                                   |
     HTTPS |                                   | HTTPS
           |                                   |
           |                                   |
           v                                   v
        JS API                               JS API
  +-----------+                             +-----------+
  |           |             Media           |           |
  Alice |  Browser  |<--------------------------->|  Browser  | Bob
  |           |           DTLS+SRTP         |           |
  +-----------+                             +-----------+
        ^      ^--+                      +--^     ^
        |         |                      |        |
        v         |                      |        v
  +-----------+   |                      |  +-----------+
  |           |<-------------------------+  |           |
  |   IdP1    |   |                         |    IdP2   |
  |           |   +------------------------>|           |
  +-----------+                             +-----------+
\end{verbatim}
\caption{A Federated Call with IdP-Based Identity}
\label{fig:federated-call-with-idp}
\end{figure}

\subsection{Initial Signaling}

For simplicity, assume the topology in Figure~\ref{fig:call-with-idp}.  Alice and Bob are
both users of a common calling service; they both have approved the
calling service to make calls (we defer the discussion of device
access permissions until later).  They are both connected to the
calling service via HTTPS and so know the origin with some level of
confidence.  They also have accounts with some IdP.  This sort of
identity service is becoming increasingly common in the Web
environment (with technologies such as Federated Google Login,
Facebook Connect, OAuth, OpenID, WebFinger), and is often provided as
a side effect service of a user's ordinary accounts with some
service.  In this example, we show Alice and Bob using a separate
identity service, though the identity service may be the same entity
as the calling service or there may be no identity service at all.

Alice is logged onto the calling service and decides to call Bob. She
can see from the calling service that he is online and the calling
service presents a JS UI in the form of a button next to Bob's name
which says ``Call''.  Alice clicks the button, which initiates a JS
callback that instantiates a PeerConnection object.  This does not
require a security check: JS from any origin is allowed to get this
far.

Once the PeerConnection is created, the calling service JS needs to
set up some media.  Because this is an audio/video call, it creates a
MediaStream with two MediaStreamTracks, one connected to an audio
input and one connected to a video input.  At this point, the first
security check is required: untrusted origins are not allowed to
access the camera and microphone, so the browser prompts Alice for
permission.

In the current W3C API, once some streams have been added, Alice's
browser + JS generates a signaling message [RFC8829] containing:

\begin{itemize}
\item Media channel information

\item Interactive Connectivity Establishment (ICE) [RFC8445] candidates

\item A ``fingerprint'' attribute binding the communication to a key pair
  [RFC5763].  Note that this key may simply be ephemerally generated
  for this call or specific to this domain, and Alice may have a
  large number of such keys.
\end{itemize}

Prior to sending out the signaling message, the PeerConnection code
contacts the identity service and obtains an assertion binding
Alice's identity to her fingerprint.  The exact details depend on the
identity service (though as discussed in Section~\ref{sec:web-based-peer-auth} PeerConnection can
be agnostic to them), but for now it's easiest to think of as an
OAuth token.  The assertion may bind other information to the
identity besides the fingerprint, but at minimum it needs to bind the
fingerprint.

This message is sent to the signaling server, e.g., by fetch()
[fetch] or by WebSockets [RFC6455], over TLS [RFC8446].  The
signaling server processes the message from Alice's browser,
determines that this is a call to Bob, and sends a signaling message
to Bob's browser (again, the format is currently undefined).  The JS
on Bob's browser processes it, and alerts Bob to the incoming call
and to Alice's identity.  In this case, Alice has provided an
identity assertion and so Bob's browser contacts Alice's IdP (again,
this is done in a generic way so the browser has no specific
knowledge of the IdP) to verify the assertion.  It is also possible
to have IdPs with which the browser has a specific trust
relationship, as described in Section~\ref{sec:trust-relationships}.  This allows the browser
to display a trusted element in the browser chrome indicating that a
call is coming in from Alice.  If Alice is in Bob's address book,
then this interface might also include her real name, a picture, etc.
The calling site will also provide some user interface element (e.g.,
a button) to allow Bob to answer the call, though this is most likely
not part of the trusted UI.

If Bob agrees, a PeerConnection is instantiated with the message from
Alice's side.  Then, a similar process occurs as on Alice's browser:
Bob's browser prompts him for device permission, the media streams
are created, and a return signaling message containing media
information, ICE candidates, and a fingerprint is sent back to Alice
via the signaling service.  If Bob has a relationship with an IdP,
the message will also come with an identity assertion.

At this point, Alice and Bob each know that the other party wants to
have a secure call with them.  Based purely on the interface provided
by the signaling server, they know that the signaling server claims
that the call is from Alice to Bob. This level of security is
provided merely by having the fingerprint in the message and having
that message received securely from the signaling server.  Because
the far end sent an identity assertion along with their message, they
know that this is verifiable from the IdP as well.  Note that if the
call is federated, as shown in Figure~\ref{fig:federated-call-with-idp}, then Alice is able to verify
Bob's identity in a way that is not mediated by either her signaling
server or Bob's.  Rather, she verifies it directly with Bob's IdP.

Of course, the call works perfectly well if either Alice or Bob
doesn't have a relationship with an IdP; they just get a lower level
of assurance.  I.e., they simply have whatever information their
calling site claims about the caller/callee's identity.  Moreover,
Alice might wish to make an anonymous call through an anonymous
calling site, in which case she would of course just not provide any
identity assertion and the calling site would mask her identity from
Bob.

\subsection{Media Consent Verification}

As described in [RFC8826], Section 4.2, media consent verification is
provided via ICE.  Thus, Alice and Bob perform ICE checks with each
other.  At the completion of these checks, they are ready to send
non-ICE data.

At this point, Alice knows that (a) Bob (assuming he is verified via
his IdP) or someone else who the signaling service is claiming is Bob
is willing to exchange traffic with her and (b) either Bob is at the
IP address which she has verified via ICE or there is an attacker who
is on-path to that IP address detouring the traffic.  Note that it is
not possible for an attacker who is on-path between Alice and Bob but
not attached to the signaling service to spoof these checks because
they do not have the ICE credentials.  Bob has the same security
guarantees with respect to Alice.

\subsection{DTLS Handshake}

Once the requisite ICE checks have completed, Alice and Bob can set
up a secure channel or channels.  This is performed via DTLS
[RFC6347] and DTLS-SRTP [RFC5763] keying for SRTP [RFC3711] for the
media channel and the Stream Control Transmission Protocol (SCTP)
over DTLS [RFC8261] for data channels.  Specifically, Alice and Bob
perform a DTLS handshake on every component which has been
established by ICE.  The total number of channels depends on the
amount of muxing; in the most likely case, we are using both RTP/RTCP
mux and muxing multiple media streams on the same channel, in which
case there is only one DTLS handshake.  Once the DTLS handshake has
completed, the keys are exported [RFC5705] and used to key SRTP for
the media channels.

At this point, Alice and Bob know that they share a set of secure
data and/or media channels with keys which are not known to any
third-party attacker.  If Alice and Bob authenticated via their IdPs,
then they also know that the signaling service is not mounting a man-in-the-middle attack on their traffic.  Even if they do not use an
IdP, as long as they have minimal trust in the signaling service not
to perform a man-in-the-middle attack, they know that their
communications are secure against the signaling service as well
(i.e., that the signaling service cannot mount a passive attack on
the communications).

\subsection{Communications and Consent Freshness}

From a security perspective, everything from here on in is a little
anticlimactic: Alice and Bob exchange data protected by the keys
negotiated by DTLS.  Because of the security guarantees discussed in
the previous sections, they know that the communications are
encrypted and authenticated.

The one remaining security property we need to establish is ``consent
freshness'', i.e., allowing Alice to verify that Bob is still prepared
to receive her communications so that Alice does not continue to send
large traffic volumes to entities which went abruptly offline.  ICE
specifies periodic Session Traversal Utilities for NAT (STUN)
keepalives but only if media is not flowing.  Because the consent
issue is more difficult here, we require WebRTC implementations to
periodically send keepalives using the consent freshness mechanism
specified in [RFC7675].  If a keepalive fails and no new ICE channels
can be established, then the session is terminated.

\section{SDP Identity Attribute}

The SDP ``identity'' attribute is a session-level attribute that is
used by an endpoint to convey its identity assertion to its peer.
The identity-assertion value is encoded as base64, as described in
Section 4 of [RFC4648].

The procedures in this section are based on the assumption that the
identity assertion of an endpoint is bound to the fingerprints of the
endpoint.  This does not preclude the definition of alternative means
of binding an assertion to the endpoint, but such means are outside
the scope of this specification.

The semantics of multiple ``identity'' attributes within an offer or
answer are undefined.  Implementations SHOULD only include a single
``identity'' attribute in an offer or answer, and Relying Parties MAY
elect to ignore all but the first ``identity'' attribute.

\begin{description}[style=nextline]
\item[Name:] identity
\item[Value:] identity-assertion
\item[Usage Level:] session
\item[Charset Dependent:] no
\item[Default Value:] N/A
\end{description}

\noindent\textbf{Syntax:}

\begin{verbatim}
 identity-assertion       = identity-assertion-value
                            *(SP identity-extension)
 identity-assertion-value = base64
 identity-extension       = extension-name [ "=" extension-value ]
 extension-name           = token
 extension-value          = 1*(%x01-09 / %x0b-0c / %x0e-3a / %x3c-ff)
                            ; byte-string from [RFC4566]

 <ALPHA and DIGIT as defined in [RFC4566]>
 <base64 as defined in [RFC4566]>
\end{verbatim}

\noindent\textbf{Example:}

\begin{verbatim}
 a=identity:\
   eyJpZHAiOnsiZG9tYWluIjoiZXhhbXBsZS5vcmciLCJwcm90b2NvbCI6ImJvZ3Vz\
   In0sImFzc2VydGlvbiI6IntcImlkZW50aXR5XCI6XCJib2JAZXhhbXBsZS5vcmdc\
   IixcImNvbnRlbnRzXCI6XCJhYmNkZWZnaGlqa2xtbm9wcXJzdHV2d3l6XCIsXCJz\
   aWduYXR1cmVcIjpcIjAxMDIwMzA0MDUwNlwifSJ9
\end{verbatim}

\begin{quote}
Note that long lines in the example are folded to meet the
column width constraints of this document; the backslash (``\textbackslash'')
at the end of a line, the carriage return that follows, and
whitespace shall be ignored.
\end{quote}

This specification does not define any extensions for the attribute.

The identity-assertion value is a JSON encoded string [RFC8259].  The
JSON object contains two keys: ``assertion'' and ``idp''.  The
``assertion'' key value contains an opaque string that is consumed by
the IdP.  The ``idp'' key value contains a dictionary with one or two
further values that identify the IdP.  See Section~\ref{sec:requesting-assertions} for more
details.

\subsection{Offer/Answer Considerations}

This section defines the SDP offer/answer [RFC3264] considerations
for the SDP ``identity'' attribute.

Within this section, `initial offer' refers to the first offer in the
SDP session that contains an SDP ``identity'' attribute.

\subsubsection{Generating the Initial SDP Offer}

When an offerer sends an offer, in order to provide its identity
assertion to the peer, it includes an ``identity'' attribute in the
offer.  In addition, the offerer includes one or more SDP
``fingerprint'' attributes.  The ``identity'' attribute MUST be bound to
all the ``fingerprint'' attributes in the session description.

\subsubsection{Generating an SDP Answer}

If the answerer elects to include an ``identity'' attribute, it follows
the same steps as those in Section 5.1.1.  The answerer can choose to
include or omit an ``identity'' attribute independently, regardless of
whether the offerer did so.

\subsubsection{Processing an SDP Offer or Answer}

When an endpoint receives an offer or answer that contains an
``identity'' attribute, the answerer can use the attribute information
to contact the IdP and verify the identity of the peer.  If the
identity requires a third-party IdP as described in Section~\ref{sec:trust-relationships}, then
that IdP will need to have been specifically configured.  If the
identity verification fails, the answerer MUST discard the offer or
answer as malformed.

\subsubsection{Modifying the Session}

When modifying a session, if the set of fingerprints is unchanged,
then the sender MAY send the same ``identity'' attribute.  In this
case, the established identity MUST be applied to existing DTLS
connections as well as new connections established using one of those
fingerprints.  Note that [RFC8829], Section 5.2.1 requires that each
media section use the same set of fingerprints.  If a new ``identity''
attribute is received, then the receiver MUST apply that identity to
all existing connections.

If the set of fingerprints changes, then the sender MUST either send
a new ``identity'' attribute or none at all.  Because a change in
fingerprints also causes a new DTLS connection to be established, the
receiver MUST discard all previously established identities.

\section{Detailed Technical Description}

\subsection{Origin and Web Security Issues}

The basic unit of permissions for WebRTC is the origin [RFC6454].
Because the security of the origin depends on being able to
authenticate content from that origin, the origin can only be
securely established if data is transferred over HTTPS [RFC2818].
Thus, clients MUST treat HTTP and HTTPS origins as different
permissions domains.  Note: This follows directly from the origin
security model and is stated here merely for clarity.

Many Web browsers currently forbid by default any active mixed
content on HTTPS pages.  That is, when JavaScript is loaded from an
HTTP origin onto an HTTPS page, an error is displayed and the HTTP
content is not executed unless the user overrides the error.  Any
browser which enforces such a policy will also not permit access to
WebRTC functionality from mixed content pages (because they never
display mixed content).  Browsers which allow active mixed content
MUST nevertheless disable WebRTC functionality in mixed content
settings.

Note that it is possible for a page which was not mixed content to
become mixed content during the duration of the call.  The major risk
here is that the newly arrived insecure JS might redirect media to a
location controlled by the attacker.  Implementations MUST either
choose to terminate the call or display a warning at that point.

Also note that the security architecture depends on the keying
material not being available to move between origins.  However, it is
assumed that the identity assertion can be passed to anyone that the
page cares to.

\subsection{Device Permissions Model}

Implementations MUST obtain explicit user consent prior to providing
access to the camera and/or microphone.  Implementations MUST at
minimum support the following two permissions models for HTTPS
origins.

\begin{itemize}
\item Requests for one-time camera/microphone access.

\item Requests for permanent access.
\end{itemize}

Because HTTP origins cannot be securely established against network
attackers, implementations MUST refuse all permissions grants for
HTTP origins.

In addition, they SHOULD support requests for access that promise
that media from this grant will be sent to a single communicating
peer (obviously there could be other requests for other peers), e.g.,
``Call customerservice@example.org''.  The semantics of this request
are that the media stream from the camera and microphone will only be
routed through a connection which has been cryptographically verified
(through the IdP mechanism or an X.509 certificate in the DTLS-SRTP
handshake) as being associated with the stated identity.  Note that
it is unlikely that browsers would have X.509 certificates, but
servers might.  Browsers servicing such requests SHOULD clearly
indicate that identity to the user when asking for permission.  The
idea behind this type of permissions is that a user might have a
fairly narrow list of peers they are willing to communicate with,
e.g., ``my mother'' rather than ``anyone on Facebook''.  Narrow
permissions grants allow the browser to do that enforcement.

\begin{description}[style=nextline]
\item[API Requirement:] The API MUST provide a mechanism for the requesting
  JS to relinquish the ability to see or modify the media (e.g., via
  MediaStream.record()).  Combined with secure authentication of the
  communicating peer, this allows a user to be sure that the calling
  site is not accessing or modifying their conversion.

\item[UI Requirement:] The UI MUST clearly indicate when the user's camera
  and microphone are in use.  This indication MUST NOT be
  suppressible by the JS and MUST clearly indicate how to terminate
  device access, and provide a UI means to immediately stop camera/
  microphone input without the JS being able to prevent it.

\item[UI Requirement:] If the UI indication of camera/microphone use is
  displayed in the browser such that minimizing the browser window
  would hide the indication, or the JS creating an overlapping
  window would hide the indication, then the browser SHOULD stop
  camera and microphone input when the indication is hidden.  (Note:
  This may not be necessary in systems that are non-windows-based
  but that have good notifications support, such as phones.)
\end{description}

\begin{itemize}
\item Browsers MUST NOT permit permanent screen or application sharing
  permissions to be installed as a response to a JS request for
  permissions.  Instead, they must require some other user action
  such as a permissions setting or an application install experience
  to grant permission to a site.

\item Browsers MUST provide a separate dialog request for screen/
  application sharing permissions even if the media request is made
  at the same time as the request for camera and microphone
  permissions.

\item The browser MUST indicate any windows which are currently being
  shared in some unambiguous way.  Windows which are not visible
  MUST NOT be shared even if the application is being shared.  If
  the screen is being shared, then that MUST be indicated.
\end{itemize}

Browsers MAY permit the formation of data channels without any direct
user approval.  Because sites can always tunnel data through the
server, further restrictions on the data channel do not provide any
additional security.  (See Section~\ref{sec:communications-consent} for a related issue.)

Implementations which support some form of direct user authentication
SHOULD also provide a policy by which a user can authorize calls only
to specific communicating peers.  Specifically, the implementation
SHOULD provide the following interfaces/controls:

\begin{itemize}
\item Allow future calls to this verified user.

\item Allow future calls to any verified user who is in my system
  address book (this only works with address book integration, of
  course).
\end{itemize}

Implementations SHOULD also provide a different user interface
indication when calls are in progress to users whose identities are
directly verifiable.  Section~\ref{sec:communications-security} provides more on this.

\subsection{Communications Consent}
\label{sec:communications-consent}

Browser client implementations of WebRTC MUST implement ICE.  Server
gateway implementations which operate only at public IP addresses
MUST implement either full ICE or ICE-Lite [RFC8445].

Browser implementations MUST verify reachability via ICE prior to
sending any non-ICE packets to a given destination.  Implementations
MUST NOT provide the ICE transaction ID to JavaScript during the
lifetime of the transaction (i.e., during the period when the ICE
stack would accept a new response for that transaction).  The JS MUST
NOT be permitted to control the local ufrag and password, though it
of course knows it.

While continuing consent is required, the ICE [RFC8445], Section 11
keepalives use STUN Binding Indications, which are one-way and
therefore not sufficient.  The current WG consensus is to use ICE
Binding Requests for continuing consent freshness.  ICE already
requires that implementations respond to such requests, so this
approach is maximally compatible.  A separate document will profile
the ICE timers to be used; see [RFC7675].

\subsection{IP Location Privacy}

A side effect of the default ICE behavior is that the peer learns
one's IP address, which leaks large amounts of location information.
This has negative privacy consequences in some circumstances.  The
API requirements in this section are intended to mitigate this issue.
Note that these requirements are not intended to protect the user's
IP address from a malicious site.  In general, the site will learn at
least a user's server-reflexive address from any HTTP transaction.
Rather, these requirements are intended to allow a site to cooperate
with the user to hide the user's IP address from the other side of
the call.  Hiding the user's IP address from the server requires some
sort of explicit privacy-preserving mechanism on the client (e.g.,
Tor Browser \url{https://www.torproject.org/projects/torbrowser.html.en})
and is out of scope for this specification.

\begin{description}[style=nextline]
\item[API Requirement:] The API MUST provide a mechanism to allow the JS to
  suppress ICE negotiation (though perhaps to allow candidate
  gathering) until the user has decided to answer the call.  (Note:
  Determining when the call has been answered is a question for the
  JS.)  This enables a user to prevent a peer from learning their IP
  address if they elect not to answer a call and also from learning
  whether the user is online.

\item[API Requirement:] The API MUST provide a mechanism for the calling
  application JS to indicate that only TURN candidates are to be
  used.  This prevents the peer from learning one's IP address at
  all.  This mechanism MUST also permit suppression of the related
  address field, since that leaks local addresses.

\item[API Requirement:] The API MUST provide a mechanism for the calling
  application to reconfigure an existing call to add non-TURN
  candidates.  Taken together, this and the previous requirement
  allow ICE negotiation to start immediately on incoming call
  notification, thus reducing post-dial delay, but also to avoid
  disclosing the user's IP address until they have decided to
  answer.  They also allow users to completely hide their IP address
  for the duration of the call.  Finally, they allow a mechanism for
  the user to optimize performance by reconfiguring to allow non-TURN candidates during an active call if the user decides they no
  longer need to hide their IP address.
\end{description}

Note that some enterprises may operate proxies and/or NATs designed
to hide internal IP addresses from the outside world.  WebRTC
provides no explicit mechanism to allow this function.  Either such
enterprises need to proxy the HTTP/HTTPS and modify the SDP and/or
the JS, or there needs to be browser support to set the ``TURN-only''
policy regardless of the site's preferences.

Note: These requirements are intended to allow sites to conceal the
user's IP address from the peer.  For guidance on concealing the
user's IP address from the calling site see [RFC8828].

\subsection{Communications Security}
\label{sec:communications-security}

Implementations MUST support SRTP [RFC3711].  Implementations MUST
support DTLS [RFC6347] and DTLS-SRTP [RFC5763] [RFC5764] for SRTP
keying.  Implementations MUST support SCTP over DTLS [RFC8261].

All media channels MUST be secured via SRTP and the Secure Real-time
Transport Control Protocol (SRTCP).  Media traffic MUST NOT be sent
over plain (unencrypted) RTP or RTCP; that is, implementations MUST
NOT negotiate cipher suites with NULL encryption modes.  DTLS-SRTP
MUST be offered for every media channel.  WebRTC implementations MUST
NOT offer SDP security descriptions [RFC4568] or select it if
offered.  An SRTP Master Key Identifier (MKI) MUST NOT be used.

All data channels MUST be secured via DTLS.

All implementations MUST support DTLS 1.2 with the
TLS\_ECDHE\_ECDSA\_WITH\_AES\_128\_GCM\_SHA256 cipher suite and the P-256
curve [FIPS186].  Earlier drafts of this specification required DTLS
1.0 with the cipher suite TLS\_ECDHE\_ECDSA\_WITH\_AES\_128\_CBC\_SHA, and
at the time of this writing some implementations do not support DTLS
1.2; endpoints which support only DTLS 1.2 might encounter
interoperability issues.  The DTLS-SRTP protection profile
SRTP\_AES128\_CM\_HMAC\_SHA1\_80 MUST be supported for SRTP.
Implementations MUST favor cipher suites which support Forward
Secrecy (FS) over non-FS cipher suites and SHOULD favor Authenticated
Encryption with Associated Data (AEAD) over non-AEAD cipher suites.
Note: the IETF is in the process of standardizing DTLS 1.3
[TLS-DTLS13].

Implementations MUST NOT implement DTLS renegotiation and MUST reject
it with a ``no\_renegotiation'' alert if offered.

Endpoints MUST NOT implement TLS False Start [RFC7918].

\begin{description}[style=nextline]
\item[API Requirement:] The API MUST generate a new authentication key pair
  for every new call by default.  This is intended to allow for
  unlinkability.

\item[API Requirement:] The API MUST provide a means to reuse a key pair
  for calls.  This can be used to enable key continuity-based
  authentication, and could be used to amortize key generation
  costs.

\item[API Requirement:] Unless the user specifically configures an external
  key pair, different key pairs MUST be used for each origin.  (This
  avoids creating a super-cookie.)

\item[API Requirement:] When DTLS-SRTP is used, the API MUST NOT permit the
  JS to obtain the negotiated keying material.  This requirement
  preserves the end-to-end security of the media.

\item[UI Requirements:] A user-oriented client MUST provide an ``inspector''
  interface which allows the user to determine the ``security
  characteristics'' of the media.
\end{description}

The following properties SHOULD be displayed ``up-front'' in the
browser chrome, i.e., without requiring the user to ask for them:

\begin{itemize}
\item A client MUST provide a user interface through which a user may
  determine the ``security characteristics'' for currently
  displayed audio and video stream(s).

\item A client MUST provide a user interface through which a user may
  determine the ``security characteristics'' for transmissions of
  their microphone audio and camera video.

\item If the far endpoint was directly verified, either via a third-party verifiable X.509 certificate or via a Web IdP mechanism
  (see Section~\ref{sec:web-based-peer-auth}), the ``security characteristics'' MUST include
  the verified information.  X.509 identities and Web IdP
  identities have similar semantics and should be displayed in a
  similar way.
\end{itemize}

The following properties are more likely to require some ``drill-down'' from the user:

\begin{itemize}
\item The ``security characteristics'' MUST indicate the cryptographic
  algorithms in use (for example, ``AES-CBC'').

\item The ``security characteristics'' MUST indicate whether FS is
  provided.

\item The ``security characteristics'' MUST include some mechanism to
  allow an out-of-band verification of the peer, such as a
  certificate fingerprint or a Short Authentication String (SAS).
  These are compared by the peers to authenticate one another.
\end{itemize}

\section{Web-Based Peer Authentication}
\label{sec:web-based-peer-auth}

NOTE: The mechanism described in this section was designed relatively
early in the RTCWEB process.  In retrospect, the WG was too
optimistic about the enthusiasm for this kind of mechanism.  At the
time of publication, it has not been widely adopted or implemented.
It appears in this document as a description of the state of the art
as of this writing.

In a number of cases, it is desirable for the endpoint (i.e., the
browser) to be able to directly identify the endpoint on the other
side without trusting the signaling service to which they are
connected.  For instance, users may be making a call via a federated
system where they wish to get direct authentication of the other
side.  Alternately, they may be making a call on a site which they
minimally trust (such as a poker site) but to someone who has an
identity on a site they do trust (such as a social network).

Recently, a number of Web-based identity technologies (OAuth,
Facebook Connect, etc.) have been developed.  While the details vary,
what these technologies share is that they have a Web-based (i.e.,
HTTP/HTTPS) IdP which attests to Alice's identity.  For instance, if
Alice has an account at example.org, Alice could use the example.org
IdP to prove to others that Alice is alice@example.org.  The
development of these technologies allows us to separate calling from
identity provision: Alice could call you on a poker site but identify
herself as alice@example.org.

Whatever the underlying technology, the general principle is that the
party which is being authenticated is NOT the signaling site but
rather the user (and their browser).  Similarly, the Relying Party is
the browser and not the signaling site.  Thus, the browser MUST
generate the input to the IdP assertion process and display the
results of the verification process to the user in a way which cannot
be imitated by the calling site.

The mechanisms defined in this document do not require the browser to
implement any particular identity protocol or to support any
particular IdP.  Instead, this document provides a generic interface
which any IdP can implement.  Thus, new IdPs and protocols can be
introduced without change to either the browser or the calling
service.  This avoids the need to make a commitment to any particular
identity protocol, although browsers may opt to directly implement
some identity protocols in order to provide superior performance or
UI properties.

\subsection{Trust Relationships: IdPs, APs, and RPs}
\label{sec:trust-relationships}

Any federated identity protocol has three major participants:

\begin{description}[style=nextline]
\item[Authenticating Party (AP):] The entity which is trying to establish
  its identity.

\item[Identity Provider (IdP):] The entity which is vouching for the AP's
  identity.

\item[Relying Party (RP):] The entity which is trying to verify the AP's
  identity.
\end{description}

The AP and the IdP have an account relationship of some kind: the AP
registers with the IdP and is able to subsequently authenticate
directly to the IdP (e.g., with a password).  This means that the
browser must somehow know which IdP(s) the user has an account
relationship with.  This can either be something that the user
configures into the browser or that is configured at the calling site
and then provided to the PeerConnection by the Web application at the
calling site.  The use case for having this information configured
into the browser is that the user may ``log into'' the browser to bind
it to some identity.  This is becoming common in new browsers.
However, it should also be possible for the IdP information to simply
be provided by the calling application.

At a high level, there are two kinds of IdPs:

\begin{description}[style=nextline]
\item[Authoritative:] IdPs which have verifiable control of some section of
  the identity space.  For instance, in the realm of email, the
  operator of ``example.com'' has complete control of the namespace
  ending in ``@example.com''.  Thus, ``alice@example.com'' is whoever
  the operator says it is.  Examples of systems with authoritative
  IdPs include DNSSEC, an identity system for SIP (see [RFC8224]),
  and Facebook Connect (Facebook identities only make sense within
  the context of the Facebook system).

\item[Third-Party:] IdPs which don't have control of their section of the
  identity space but instead verify users' identities via some
  unspecified mechanism and then attest to it.  Because the IdP
  doesn't actually control the namespace, RPs need to trust that the
  IdP is correctly verifying AP identities, and there can
  potentially be multiple IdPs attesting to the same section of the
  identity space.  Probably the best-known example of a third-party
  IdP is SSL/TLS certificates, where there are a large number of
  certificate authorities (CAs) all of whom can attest to any domain
  name.
\end{description}

If an AP is authenticating via an authoritative IdP, then the RP does
not need to explicitly configure trust in the IdP at all.  The
identity mechanism can directly verify that the IdP indeed made the
relevant identity assertion (a function provided by the mechanisms in
this document), and any assertion it makes about an identity for
which it is authoritative is directly verifiable.  Note that this
does not mean that the IdP might not lie, but that is a
trustworthiness judgement that the user can make at the time they
look at the identity.

By contrast, if an AP is authenticating via a third-party IdP, the RP
needs to explicitly trust that IdP (hence the need for an explicit
trust anchor list in PKI-based SSL/TLS clients).  The list of
trustable IdPs needs to be configured directly into the browser,
either by the user or potentially by the browser manufacturer.  This
is a significant advantage of authoritative IdPs and implies that if
third-party IdPs are to be supported, the potential number needs to
be fairly small.

\subsection{Overview of Operation}

In order to provide security without trusting the calling site, the
PeerConnection component of the browser must interact directly with
the IdP.  The details of the mechanism are described in the W3C API
specification, but the general idea is that the PeerConnection
component downloads JS from a specific location on the IdP dictated
by the IdP domain name.  That JS (the ``IdP proxy'') runs in an
isolated security context within the browser, and the PeerConnection
talks to it via a secure message passing channel.

Note that there are two logically separate functions here:

\begin{itemize}
\item Identity assertion generation.

\item Identity assertion verification.
\end{itemize}

The same IdP JS ``endpoint'' is used for both functions, but of course
a given IdP might behave differently and load new JS to perform one
function or the other.

\begin{figure}[ht]
\begin{verbatim}
   +--------------------------------------+
   | Browser                              |
   |                                      |
   | +----------------------------------+ |
   | | https://calling-site.example.com | |
   | |                                  | |
   | |        Calling JS Code           | |
   | |               ^                  | |
   | +---------------|------------------+ |
   |                 | API Calls          |
   |                 v                    |
   |          PeerConnection              |
   |                 ^                    |
   |                 | API Calls          |
   |     +-----------|-------------+      |   +---------------+
   |     |           v             |      |   |               |
   |     |       IdP Proxy         |<-------->|   Identity    |
   |     |                         |      |   |   Provider    |
   |     | https://idp.example.org |      |   |               |
   |     +-------------------------+      |   +---------------+
   |                                      |
   +--------------------------------------+
\end{verbatim}
\caption{Browser IdP Proxy Architecture}
\label{fig:idp-proxy}
\end{figure}

When the PeerConnection object wants to interact with the IdP, the
sequence of events is as follows:

\begin{enumerate}
\item The browser (the PeerConnection component) instantiates an IdP
  proxy.  This allows the IdP to load whatever JS is necessary into
  the proxy.  The resulting code runs in the IdP's security
  context.

\item The IdP registers an object with the browser that conforms to the
  API defined in [webrtc-api].

\item The browser invokes methods on the object registered by the IdP
  proxy to create or verify identity assertions.
\end{enumerate}

This approach allows us to decouple the browser from any particular
IdP; the browser need only know how to load the IdP's JavaScript --
the location of which is determined based on the IdP's identity --
and to call the generic API for requesting and verifying identity
assertions.  The IdP provides whatever logic is necessary to bridge
the generic protocol to the IdP's specific requirements.  Thus, a
single browser can support any number of identity protocols,
including being forward compatible with IdPs which did not exist at
the time the browser was written.

\subsection{Items for Standardization}

There are two parts to this work:

\begin{itemize}
\item The precise information from the signaling message that must be
  cryptographically bound to the user's identity and a mechanism for
  carrying assertions in JavaScript Session Establishment Protocol
  (JSEP) messages.  This is specified in Section~\ref{sec:binding-identity}.

\item The interface to the IdP, which is defined in the companion W3C
  WebRTC API specification [webrtc-api].
\end{itemize}

The WebRTC API specification also defines JavaScript interfaces that
the calling application can use to specify which IdP to use.  That
API also provides access to the assertion-generation capability and
the status of the validation process.

\subsection{Binding Identity Assertions to JSEP Offer/Answer Transactions}
\label{sec:binding-identity}

An identity assertion binds the user's identity (as asserted by the
IdP) to the SDP offer/answer exchange and specifically to the media.
In order to achieve this, the PeerConnection must provide the DTLS-SRTP fingerprint to be bound to the identity.  This is provided as a
JavaScript object (also known as a dictionary or hash) with a single
``fingerprint'' key, as shown below:

\begin{verbatim}
{
  "fingerprint":
    [
      { "algorithm": "sha-256",
        "digest": "4A:AD:B9:B1:3F:...:E5:7C:AB" },
      { "algorithm": "sha-1",
        "digest": "74:E9:76:C8:19:...:F4:45:6B" }
    ]
}
\end{verbatim}

The ``fingerprint'' value is an array of objects.  Each object in the
array contains ``algorithm'' and ``digest'' values, which correspond
directly to the algorithm and digest values in the ``fingerprint''
attribute of the SDP [RFC8122].

This object is encoded in a JSON [RFC8259] string for passing to the
IdP.  The identity assertion returned by the IdP, which is encoded in
the ``identity'' attribute, is a JSON object that is encoded as
described in Section~\ref{sec:carrying-identity}.

This structure does not need to be interpreted by the IdP or the IdP
proxy.  It is consumed solely by the RP's browser.  The IdP merely
treats it as an opaque value to be attested to.  Thus, new parameters
can be added to the assertion without modifying the IdP.

\subsubsection{Carrying Identity Assertions}
\label{sec:carrying-identity}

Once an IdP has generated an assertion (see Section~\ref{sec:requesting-assertions}), it is
attached to the SDP offer/answer message.  This is done by adding a
new ``identity'' attribute to the SDP.  The sole contents of this value
is the identity assertion.  The identity assertion produced by the
IdP is encoded into a UTF-8 JSON text, then base64-encoded [RFC4648]
to produce this string.  For example:

\begin{verbatim}
v=0
o=- 1181923068 1181923196 IN IP4 ua1.example.com
s=example1
c=IN IP4 ua1.example.com
a=fingerprint:sha-1 \
  4A:AD:B9:B1:3F:82:18:3B:54:02:12:DF:3E:5D:49:6B:19:E5:7C:AB
a=identity:\
  eyJpZHAiOnsiZG9tYWluIjoiZXhhbXBsZS5vcmciLCJwcm90b2NvbCI6ImJvZ3Vz\
  In0sImFzc2VydGlvbiI6IntcImlkZW50aXR5XCI6XCJib2JAZXhhbXBsZS5vcmdc\
  IixcImNvbnRlbnRzXCI6XCJhYmNkZWZnaGlqa2xtbm9wcXJzdHV2d3l6XCIsXCJz\
  aWduYXR1cmVcIjpcIjAxMDIwMzA0MDUwNlwifSJ9
a=...
t=0 0
m=audio 6056 RTP/SAVP 0
a=sendrecv
...
\end{verbatim}

\begin{quote}
Note that long lines in the example are folded to meet the
column width constraints of this document; the backslash (``\textbackslash'')
at the end of a line, the carriage return that follows, and
whitespace shall be ignored.
\end{quote}

The ``identity'' attribute attests to all ``fingerprint'' attributes in
the session description.  It is therefore a session-level attribute.

Multiple ``fingerprint'' values can be used to offer alternative
certificates for a peer.  The ``identity'' attribute MUST include all
``fingerprint'' values that are included in ``fingerprint'' attributes of
the session description.

The RP browser MUST verify that the in-use certificate for a DTLS
connection is in the set of fingerprints returned from the IdP when
verifying an assertion.

\subsection{Determining the IdP URI}

In order to ensure that the IdP is under control of the domain owner
rather than someone who merely has an account on the domain owner's
server (e.g., in shared hosting scenarios), the IdP JavaScript is
hosted at a deterministic location based on the IdP's domain name.
Each IdP proxy instance is associated with two values:

\begin{description}[style=nextline]
\item[authority:] The authority [RFC3986] at which the IdP's service is
  hosted.

\item[protocol:] The specific IdP protocol which the IdP is using.  This is
  a completely opaque IdP-specific string, but allows an IdP to
  implement two protocols in parallel.  This value may be the empty
  string.  If no value for protocol is provided, a value of
  ``default'' is used.
\end{description}

Each IdP MUST serve its initial entry page (i.e., the one loaded by
the IdP proxy) from a well-known URI [RFC8615].  The well-known URI
for an IdP proxy is formed from the following URI components:

\begin{enumerate}
\item The scheme, ``https:''.  An IdP MUST be loaded using HTTPS
  [RFC2818].

\item The authority [RFC3986].  As noted above, the authority MAY
  contain a non-default port number or userinfo sub-component.
  Both are removed when determining if an asserted identity matches
  the name of the IdP.

\item The path, starting with ``/.well-known/idp-proxy/'' and appended
  with the IdP protocol.  Note that the separator characters `/'
  (\%2F) and `\textbackslash' (\%5C) MUST NOT be permitted in the protocol field,
  lest an attacker be able to direct requests outside of the
  controlled ``/.well-known/'' prefix.  Query and fragment values MAY
  be used by including `?' or `\#' characters.
\end{enumerate}

For example, for the IdP ``identity.example.com'' and the protocol
``example'', the URL would be:

\url{https://identity.example.com/.well-known/idp-proxy/example}

The IdP MAY redirect requests to this URL, but they MUST retain the
``https:'' scheme.  This changes the effective origin of the IdP, but
not the domain of the identities that the IdP is permitted to assert
and validate.  I.e., the IdP is still regarded as authoritative for
the original domain.

\subsubsection{Authenticating Party}

How an AP determines the appropriate IdP domain is out of scope of
this specification.  In general, however, the AP has some actual
account relationship with the IdP, as this identity is what the IdP
is attesting to.  Thus, the AP somehow supplies the IdP information
to the browser.  Some potential mechanisms include:

\begin{itemize}
\item Provided by the user directly.

\item Selected from some set of IdPs known to the calling site (e.g., a
  button that shows ``Authenticate via Facebook Connect'').
\end{itemize}

\subsubsection{Relying Party}

Unlike the AP, the RP need not have any particular relationship with
the IdP.  Rather, it needs to be able to process whatever assertion
is provided by the AP.  As the assertion contains the IdP's identity
in the ``idp'' field of the JSON-encoded object (see Section~\ref{sec:requesting-assertions}), the
URI can be constructed directly from the assertion, and thus the RP
can directly verify the technical validity of the assertion with no
user interaction.  Authoritative assertions need only be verifiable.
Third-party assertions also MUST be verified against local policy, as
described in Section~\ref{sec:identity-formats}.

\subsection{Requesting Assertions}
\label{sec:requesting-assertions}

The input to the identity assertion generation process is the JSON-encoded object described in Section~\ref{sec:binding-identity} that contains the set of
certificate fingerprints the browser intends to use.  This string is
treated as opaque from the perspective of the IdP.

The browser also identifies the origin that the PeerConnection is run
in, which allows the IdP to make decisions based on who is requesting
the assertion.

An application can optionally provide a user identifier hint when
specifying an IdP.  This value is a hint that the IdP can use to
select amongst multiple identities, or to avoid providing assertions
for unwanted identities.  The ``username'' is a string that has no
meaning to any entity other than the IdP; it can contain any data the
IdP needs in order to correctly generate an assertion.

An identity assertion that is successfully provided by the IdP
consists of the following information:

\begin{description}[style=nextline]
\item[idp:] The domain name of an IdP and the protocol string.  This MAY
  identify a different IdP or protocol from the one that generated
  the assertion.

\item[assertion:] An opaque value containing the assertion itself.  This is
  only interpretable by the identified IdP or the IdP code running
  in the client.
\end{description}

Figure~\ref{fig:example-assertion} shows an example assertion formatted as JSON.  In this case,
the message has presumably been digitally signed/MACed in some way
that the IdP can later verify it, but this is an implementation
detail and out of scope of this document.

\begin{figure}[ht]
\begin{verbatim}
{
  "idp":{
    "domain": "example.org",
    "protocol": "bogus"
  },
  "assertion": "{\"identity\":\"bob@example.org\",
                 \"contents\":\"abcdefghijklmnopqrstuvwyz\",
                 \"signature\":\"010203040506\"}"
}
\end{verbatim}
\caption{Example Assertion}
\label{fig:example-assertion}
\end{figure}

For use in signaling, the assertion is serialized into JSON,
base64-encoded [RFC4648], and used as the value of the ``identity''
attribute.  IdPs SHOULD ensure that any assertions they generate
cannot be interpreted in a different context.  E.g., they should use
a distinct format or have separate cryptographic keys for assertion
generation and other purposes.  Line breaks are inserted solely for
readability.

\subsection{Managing User Login}

In order to generate an identity assertion, the IdP needs proof of
the user's identity.  It is common practice to authenticate users
(using passwords or multi-factor authentication), then use cookies
[RFC6265] or HTTP authentication [RFC7617] for subsequent exchanges.

The IdP proxy is able to access cookies, HTTP authentication data, or
other persistent session data because it operates in the security
context of the IdP origin.  Therefore, if a user is logged in, the
IdP could have all the information needed to generate an assertion.

An IdP proxy is unable to generate an assertion if the user is not
logged in, or the IdP wants to interact with the user to acquire more
information before generating the assertion.  If the IdP wants to
interact with the user before generating an assertion, the IdP proxy
can fail to generate an assertion and instead indicate a URL where
login should proceed.

The application can then load the provided URL to enable the user to
enter credentials.  The communication between the application and the
IdP is described in [webrtc-api].

\section{Verifying Assertions}

The input to identity validation is the assertion string taken from a
decoded ``identity'' attribute.

The IdP proxy verifies the assertion.  Depending on the identity
protocol, the proxy might contact the IdP server or other servers.
For instance, an OAuth-based protocol will likely require using the
IdP as an oracle, whereas with a signature-based scheme it might be
able to verify the assertion without contacting the IdP, provided
that it has cached the relevant public key.

Regardless of the mechanism, if verification succeeds, a successful
response from the IdP proxy consists of the following information:

\begin{description}[style=nextline]
\item[identity:] The identity of the AP from the IdP's perspective.
  Details of this are provided in Section~\ref{sec:identity-formats}.

\item[contents:] The original unmodified string provided by the AP as input
  to the assertion generation process.
\end{description}

Figure~\ref{fig:example-verification} shows an example response, which is JSON-formatted.

\begin{figure}[ht]
\begin{verbatim}
{
  "identity": "bob@example.org",
  "contents": "{\"fingerprint\":[ ... ]}"
}
\end{verbatim}
\caption{Example Verification Result}
\label{fig:example-verification}
\end{figure}

\subsection{Identity Formats}
\label{sec:identity-formats}

The identity provided from the IdP to the RP browser MUST consist of
a string representing the user's identity.  This string is in the
form ``\textless user\textgreater @\textless domain\textgreater'', where ``user'' consists of any character, and
domain is an internationalized domain name [RFC5890] encoded as a
sequence of U-labels.

The PeerConnection API MUST check this string as follows:

\begin{enumerate}
\item If the ``domain'' portion of the string is equal to the domain name
  of the IdP proxy, then the assertion is valid, as the IdP is
  authoritative for this domain.  Comparison of domain names is
  done using the label equivalence rule defined in Section 2.3.2.4
  of [RFC5890].

\item If the ``domain'' portion of the string is not equal to the domain
  name of the IdP proxy, then the PeerConnection object MUST reject
  the assertion unless both:

  \begin{enumerate}
  \item the IdP domain is trusted as an acceptable third-party IdP;
    and

  \item local policy is configured to trust this IdP domain for the
    domain portion of the identity string.
  \end{enumerate}
\end{enumerate}

Any `@' or `\%' characters in the ``user'' portion of the identity MUST
be escaped according to the ``percent-encoding'' rules defined in
Section 2.1 of [RFC3986].  Characters other than `@' and `\%' MUST NOT
be percent-encoded.  For example, with a ``user'' of ``user@133'' and a
``domain'' of ``identity.example.com'', the resulting string will be
encoded as ``user\%40133@identity.example.com''.

Implementations are cautioned to take care when displaying user
identities containing escaped `@' characters.  If such characters are
unescaped prior to display, implementations MUST distinguish between
the domain of the IdP proxy and any domain that might be implied by
the portion of the ``\textless user\textgreater'' portion that appears after the escaped
``@'' sign.

\section{Security Considerations}

Much of the security analysis of RTCWEB is contained in [RFC8826] or
in the discussion of the particular issues above.  In order to avoid
repetition, this section focuses on (a) residual threats that are not
addressed by this document and (b) threats produced by failure/misbehavior of one of the components in the system.

\subsection{Communications Security}

If HTTPS is not used to secure communications to the signaling
server, and the identity mechanism used in Section~\ref{sec:web-based-peer-auth} is not used,
then any on-path attacker can replace the DTLS-SRTP fingerprints in
the handshake and thus substitute its own identity for that of either
endpoint.

Even if HTTPS is used, the signaling server can potentially mount a
man-in-the-middle attack unless implementations have some mechanism
for independently verifying keys.  The UI requirements in Section~\ref{sec:communications-security}
are designed to provide such a mechanism for motivated/security
conscious users, but are not suitable for general use.  The identity
service mechanisms in Section~\ref{sec:web-based-peer-auth} are more suitable for general use.
Note, however, that a malicious signaling service can strip off any
such identity assertions, though it cannot forge new ones.  Note that
all of the third-party security mechanisms available (whether X.509
certificates or a third-party IdP) rely on the security of the third
party -- this is of course also true of the user's connection to the
Web site itself.  Users who wish to assure themselves of security
against a malicious IdP can only do so by verifying peer credentials
directly, e.g., by checking the peer's fingerprint against a value
delivered out of band.

In order to protect against malicious content JavaScript, that
JavaScript MUST NOT be allowed to have direct access to -- or perform
computations with -- DTLS keys.  For instance, if content JS were
able to compute digital signatures, then it would be possible for
content JS to get an identity assertion for a browser's generated key
and then use that assertion plus a signature by the key to
authenticate a call protected under an ephemeral Diffie-Hellman (DH)
key controlled by the content JS, thus violating the security
guarantees otherwise provided by the IdP mechanism.  Note that it is
not sufficient merely to deny the content JS direct access to the
keys, as some have suggested doing with the WebCrypto API
[webcrypto].  The JS must also not be allowed to perform operations
that would be valid for a DTLS endpoint.  By far the safest approach
is simply to deny the ability to perform any operations that depend
on secret information associated with the key.  Operations that
depend on public information, such as exporting the public key, are
of course safe.

\subsection{Privacy}

The requirements in this document are intended to allow:

\begin{itemize}
\item Users to participate in calls without revealing their location.

\item Potential callees to avoid revealing their location and even
  presence status prior to agreeing to answer a call.
\end{itemize}

However, these privacy protections come at a performance cost in
terms of using TURN relays and, in the latter case, delaying ICE.
Sites SHOULD make users aware of these tradeoffs.

Note that the protections provided here assume a non-malicious
calling service.  As the calling service always knows the user's
status and (absent the use of a technology like Tor) their IP
address, they can violate the user's privacy at will.  Users who wish
privacy against the calling sites they are using must use separate
privacy-enhancing technologies such as Tor. Combined WebRTC/Tor
implementations SHOULD arrange to route the media as well as the
signaling through Tor. Currently this will produce very suboptimal
performance.

Additionally, any identifier which persists across multiple calls is
potentially a problem for privacy, especially for anonymous calling
services.  Such services SHOULD instruct the browser to use separate
DTLS keys for each call and also to use TURN throughout the call.
Otherwise, the other side will learn linkable information that would
allow them to correlate the browser across multiple calls.
Additionally, browsers SHOULD implement the privacy-preserving CNAME
generation mode of [RFC7022].

\subsection{Denial of Service}

The consent mechanisms described in this document are intended to
mitigate denial-of-service (DoS) attacks in which an attacker uses
clients to send large amounts of traffic to a victim without the
consent of the victim.  While these mechanisms are sufficient to
protect victims who have not implemented WebRTC at all, WebRTC
implementations need to be more careful.

Consider the case of a call center which accepts calls via WebRTC.
An attacker proxies the call center's front-end and arranges for
multiple clients to initiate calls to the call center.  Note that
this requires user consent in many cases, but because the data
channel does not need consent, they can use that directly.  Since ICE
will complete, browsers can then be induced to send large amounts of
data to the victim call center if it supports the data channel at
all.  Preventing this attack requires that automated WebRTC
implementations implement sensible flow control and have the ability
to triage out (i.e., stop responding to ICE probes on) calls which
are behaving badly, and especially to be prepared to remotely
throttle the data channel in the absence of plausible audio and video
(which the attacker cannot control).

Another related attack is for the signaling service to swap the ICE
candidates for the audio and video streams, thus forcing a browser to
send video to the sink that the other victim expects will contain
audio (perhaps it is only expecting audio!), potentially causing
overload.  Muxing multiple media flows over a single transport makes
it harder to individually suppress a single flow by denying ICE
keepalives.  Either media-level (RTCP) mechanisms must be used or the
implementation must deny responses entirely, thus terminating the
call.

Yet another attack, suggested by Magnus Westerlund, is for the
attacker to cross-connect offers and answers as follows.  It induces
the victim to make a call and then uses its control of other users'
browsers to get them to attempt a call to someone.  It then
translates their offers into apparent answers to the victim, which
looks like large-scale parallel forking.  The victim still responds
to ICE responses, and now the browsers all try to send media to the
victim.  Implementations can defend themselves from this attack by
only responding to ICE Binding Requests for a limited number of
remote ufrags (this is the reason for the requirement that the JS not
be able to control the ufrag and password).  [RFC8834], Section 13
documents a number of potential RTCP-based DoS attacks and
countermeasures.

Note that attacks based on confusing one end or the other about
consent are possible even in the face of the third-party identity
mechanism as long as major parts of the signaling messages are not
signed.  On the other hand, signing the entire message severely
restricts the capabilities of the calling application, so there are
difficult tradeoffs here.

\subsection{IdP Authentication Mechanism}

This mechanism relies for its security on the IdP and on the
PeerConnection correctly enforcing the security invariants described
above.  At a high level, the IdP is attesting that the user
identified in the assertion wishes to be associated with the
assertion.  Thus, it must not be possible for arbitrary third parties
to get assertions tied to a user or to produce assertions that RPs
will accept.

\subsubsection{PeerConnection Origin Check}

Fundamentally, the IdP proxy is just a piece of HTML and JS loaded by
the browser, so nothing stops a Web attacker from creating their own
IFRAME, loading the IdP proxy HTML/JS, and requesting a signature
over their own keys rather than those generated in the browser.
However, that proxy would be in the attacker's origin, not the IdP's
origin.  Only the browser itself can instantiate a context that
(a) is in the IdP's origin and (b) exposes the correct API surface.
Thus, the IdP proxy on the sender's side MUST ensure that it is
running in the IdP's origin prior to issuing assertions.

Note that this check only asserts that the browser (or some other
entity with access to the user's authentication data) attests to the
request and hence to the fingerprint.  It does not demonstrate that
the browser has access to the associated private key, and therefore
an attacker can attach their own identity to another party's keying
material, thus making a call which comes from Alice appear to come
from the attacker.  See [RFC8844] for defenses against this form of
attack.

\subsubsection{IdP Well-Known URI}

As described in Section 7.5, the IdP proxy HTML/JS landing page is
located at a well-known URI based on the IdP's domain name.  This
requirement prevents an attacker who can write some resources at the
IdP (e.g., on one's Facebook wall) from being able to impersonate the
IdP.

\subsubsection{Privacy of IdP-Generated Identities and the Hosting Site}

Depending on the structure of the IdP's assertions, the calling site
may learn the user's identity from the perspective of the IdP.  In
many cases, this is not an issue because the user is authenticating
to the site via the IdP in any case -- for instance, when the user
has logged in with Facebook Connect and is then authenticating their
call with a Facebook identity.  However, in other cases, the user may
not have already revealed their identity to the site.  In general,
IdPs SHOULD either verify that the user is willing to have their
identity revealed to the site (e.g., through the usual IdP
permissions dialog) or arrange that the identity information is only
available to known RPs (e.g., social graph adjacencies) but not to
the calling site.  The ``domain'' field of the assertion request can be
used to check that the user has agreed to disclose their identity to
the calling site; because it is supplied by the PeerConnection it can
be trusted to be correct.

\subsubsection{Security of Third-Party IdPs}

As discussed above, each third-party IdP represents a new universal
trust point and therefore the number of these IdPs needs to be quite
limited.  Most IdPs, even those which issue unqualified identities
such as Facebook, can be recast as authoritative IdPs (e.g.,
123456@facebook.com).  However, in such cases, the user interface
implications are not entirely desirable.  One intermediate approach
is to have special (potentially user configurable) UI for large
authoritative IdPs, thus allowing the user to instantly grasp that
the call is being authenticated by Facebook, Google, etc.

\paragraph{Confusable Characters}

Because a broad range of characters are permitted in identity
strings, it may be possible for attackers to craft identities which
are confusable with other identities (see [RFC6943] for more on this
topic).  This is a problem with any identifier space of this type
(e.g., email addresses).  Those minting identifiers should avoid
mixed scripts and similar confusable characters.  Those presenting
these identifiers to a user should consider highlighting cases of
mixed script usage (see [RFC5890], Section 4.4).  Other best
practices are still in development.

\subsubsection{Web Security Feature Interactions}

A number of optional Web security features have the potential to
cause issues for this mechanism, as discussed below.

\paragraph{Popup Blocking}

When popup blocking is in use, the IdP proxy is unable to generate
popup windows, dialogs, or any other form of user interactions.  This
prevents the IdP proxy from being used to circumvent user
interaction.  The ``LOGINNEEDED'' message allows the IdP proxy to
inform the calling site of a need for user login, providing the
information necessary to satisfy this requirement without resorting
to direct user interaction from the IdP proxy itself.

\paragraph{Third Party Cookies}

Some browsers allow users to block third party cookies (cookies
associated with origins other than the top-level page) for privacy
reasons.  Any IdP which uses cookies to persist logins will be broken
by third-party cookie blocking.  One option is to accept this as a
limitation; another is to have the PeerConnection object disable
third-party cookie blocking for the IdP proxy.

\section{IANA Considerations}

This specification defines the ``identity'' SDP attribute per the
procedures of Section 8.2.4 of [RFC4566].  The required information
for the registration is included here:

\begin{description}[style=nextline]
\item[Contact Name:] IESG (iesg@ietf.org)
\item[Attribute Name:] identity
\item[Long Form:] identity
\item[Type of Attribute:] session
\item[Charset Considerations:] This attribute is not subject to the charset
  attribute.
\item[Purpose:] This attribute carries an identity assertion, binding an
  identity to the transport-level security session.
\item[Appropriate Values:] See Section 5 of RFC 8827.
\item[Mux Category:] NORMAL
\end{description}

This section registers the ``idp-proxy'' well-known URI from [RFC8615].

\begin{description}[style=nextline]
\item[URI suffix:] idp-proxy
\item[Change controller:] IETF
\end{description}

\section{References}

\subsection{Normative References}

\begin{description}[style=nextline, leftmargin=2cm, labelwidth=1.8cm]

\item[{[FIPS186]}] National Institute of Standards and Technology (NIST),
  ``Digital Signature Standard (DSS)'', NIST PUB 186-4,
  DOI 10.6028/NIST.FIPS.186-4, July 2013,
  \url{https://doi.org/10.6028/NIST.FIPS.186-4}.

\item[{[RFC2119]}] Bradner, S., ``Key words for use in RFCs to Indicate
  Requirement Levels'', BCP 14, RFC 2119,
  DOI 10.17487/RFC2119, March 1997,
  \url{https://www.rfc-editor.org/info/rfc2119}.

\item[{[RFC2818]}] Rescorla, E., ``HTTP Over TLS'', RFC 2818,
  DOI 10.17487/RFC2818, May 2000,
  \url{https://www.rfc-editor.org/info/rfc2818}.

\item[{[RFC3264]}] Rosenberg, J. and H. Schulzrinne, ``An Offer/Answer Model
  with Session Description Protocol (SDP)'', RFC 3264,
  DOI 10.17487/RFC3264, June 2002,
  \url{https://www.rfc-editor.org/info/rfc3264}.

\item[{[RFC3711]}] Baugher, M., McGrew, D., Naslund, M., Carrara, E., and K.
  Norrman, ``The Secure Real-time Transport Protocol (SRTP)'',
  RFC 3711, DOI 10.17487/RFC3711, March 2004,
  \url{https://www.rfc-editor.org/info/rfc3711}.

\item[{[RFC3986]}] Berners-Lee, T., Fielding, R., and L. Masinter, ``Uniform
  Resource Identifier (URI): Generic Syntax'', STD 66,
  RFC 3986, DOI 10.17487/RFC3986, January 2005,
  \url{https://www.rfc-editor.org/info/rfc3986}.

\item[{[RFC4566]}] Handley, M., Jacobson, V., and C. Perkins, ``SDP: Session
  Description Protocol'', RFC 4566, DOI 10.17487/RFC4566,
  July 2006, \url{https://www.rfc-editor.org/info/rfc4566}.

\item[{[RFC4568]}] Andreasen, F., Baugher, M., and D. Wing, ``Session
  Description Protocol (SDP) Security Descriptions for Media
  Streams'', RFC 4568, DOI 10.17487/RFC4568, July 2006,
  \url{https://www.rfc-editor.org/info/rfc4568}.

\item[{[RFC4648]}] Josefsson, S., ``The Base16, Base32, and Base64 Data
  Encodings'', RFC 4648, DOI 10.17487/RFC4648, October 2006,
  \url{https://www.rfc-editor.org/info/rfc4648}.

\item[{[RFC5763]}] Fischl, J., Tschofenig, H., and E. Rescorla, ``Framework
  for Establishing a Secure Real-time Transport Protocol
  (SRTP) Security Context Using Datagram Transport Layer
  Security (DTLS)'', RFC 5763, DOI 10.17487/RFC5763, May
  2010, \url{https://www.rfc-editor.org/info/rfc5763}.

\item[{[RFC5764]}] McGrew, D. and E. Rescorla, ``Datagram Transport Layer
  Security (DTLS) Extension to Establish Keys for the Secure
  Real-time Transport Protocol (SRTP)'', RFC 5764,
  DOI 10.17487/RFC5764, May 2010,
  \url{https://www.rfc-editor.org/info/rfc5764}.

\item[{[RFC5890]}] Klensin, J., ``Internationalized Domain Names for
  Applications (IDNA): Definitions and Document Framework'',
  RFC 5890, DOI 10.17487/RFC5890, August 2010,
  \url{https://www.rfc-editor.org/info/rfc5890}.

\item[{[RFC6347]}] Rescorla, E. and N. Modadugu, ``Datagram Transport Layer
  Security Version 1.2'', RFC 6347, DOI 10.17487/RFC6347,
  January 2012, \url{https://www.rfc-editor.org/info/rfc6347}.

\item[{[RFC6454]}] Barth, A., ``The Web Origin Concept'', RFC 6454,
  DOI 10.17487/RFC6454, December 2011,
  \url{https://www.rfc-editor.org/info/rfc6454}.

\item[{[RFC7022]}] Begen, A., Perkins, C., Wing, D., and E. Rescorla,
  ``Guidelines for Choosing RTP Control Protocol (RTCP)
  Canonical Names (CNAMEs)'', RFC 7022, DOI 10.17487/RFC7022,
  September 2013, \url{https://www.rfc-editor.org/info/rfc7022}.

\item[{[RFC7675]}] Perumal, M., Wing, D., Ravindranath, R., Reddy, T., and M.
  Thomson, ``Session Traversal Utilities for NAT (STUN) Usage
  for Consent Freshness'', RFC 7675, DOI 10.17487/RFC7675,
  October 2015, \url{https://www.rfc-editor.org/info/rfc7675}.

\item[{[RFC7918]}] Langley, A., Modadugu, N., and B. Moeller, ``Transport
  Layer Security (TLS) False Start'', RFC 7918,
  DOI 10.17487/RFC7918, August 2016,
  \url{https://www.rfc-editor.org/info/rfc7918}.

\item[{[RFC8122]}] Lennox, J. and C. Holmberg, ``Connection-Oriented Media
  Transport over the Transport Layer Security (TLS) Protocol
  in the Session Description Protocol (SDP)'', RFC 8122,
  DOI 10.17487/RFC8122, March 2017,
  \url{https://www.rfc-editor.org/info/rfc8122}.

\item[{[RFC8174]}] Leiba, B., ``Ambiguity of Uppercase vs Lowercase in RFC
  2119 Key Words'', BCP 14, RFC 8174, DOI 10.17487/RFC8174,
  May 2017, \url{https://www.rfc-editor.org/info/rfc8174}.

\item[{[RFC8259]}] Bray, T., Ed., ``The JavaScript Object Notation (JSON) Data
  Interchange Format'', STD 90, RFC 8259,
  DOI 10.17487/RFC8259, December 2017,
  \url{https://www.rfc-editor.org/info/rfc8259}.

\item[{[RFC8261]}] Tuexen, M., Stewart, R., Jesup, R., and S. Loreto,
  ``Datagram Transport Layer Security (DTLS) Encapsulation of
  SCTP Packets'', RFC 8261, DOI 10.17487/RFC8261, November
  2017, \url{https://www.rfc-editor.org/info/rfc8261}.

\item[{[RFC8445]}] Keranen, A., Holmberg, C., and J. Rosenberg, ``Interactive
  Connectivity Establishment (ICE): A Protocol for Network
  Address Translator (NAT) Traversal'', RFC 8445,
  DOI 10.17487/RFC8445, July 2018,
  \url{https://www.rfc-editor.org/info/rfc8445}.

\item[{[RFC8446]}] Rescorla, E., ``The Transport Layer Security (TLS) Protocol
  Version 1.3'', RFC 8446, DOI 10.17487/RFC8446, August 2018,
  \url{https://www.rfc-editor.org/info/rfc8446}.

\item[{[RFC8615]}] Nottingham, M., ``Well-Known Uniform Resource Identifiers
  (URIs)'', RFC 8615, DOI 10.17487/RFC8615, May 2019,
  \url{https://www.rfc-editor.org/info/rfc8615}.

\item[{[RFC8825]}] Alvestrand, H., ``Overview: Real-Time Protocols for
  Browser-Based Applications'', RFC 8825,
  DOI 10.17487/RFC8825, January 2021,
  \url{https://www.rfc-editor.org/info/rfc8825}.

\item[{[RFC8826]}] Rescorla, E., ``Security Considerations for WebRTC'',
  RFC 8826, DOI 10.17487/RFC8826, January 2021,
  \url{https://www.rfc-editor.org/info/rfc8826}.

\item[{[RFC8829]}] Uberti, J., Jennings, C., and E. Rescorla, Ed.,
  ``JavaScript Session Establishment Protocol (JSEP)'',
  RFC 8829, DOI 10.17487/RFC8829, January 2021,
  \url{https://www.rfc-editor.org/info/rfc8829}.

\item[{[RFC8834]}] Perkins, C., Westerlund, M., and J. Ott, ``Media Transport
  and Use of RTP in WebRTC'', RFC 8834, DOI 10.17487/RFC8834,
  January 2021, \url{https://www.rfc-editor.org/info/rfc8834}.

\item[{[RFC8844]}] Thomson, M. and E. Rescorla, ``Unknown Key-Share Attacks on
  Uses of TLS with the Session Description Protocol (SDP)'',
  RFC 8844, DOI 10.17487/RFC8844, January 2021,
  \url{https://www.rfc-editor.org/info/rfc8844}.

\item[{[webcrypto]}] Watson, M., ``Web Cryptography API'', W3C Recommendation, 26
  January 2017,
  \url{https://www.w3.org/TR/2017/REC-WebCryptoAPI-20170126/}.

\item[{[webrtc-api]}] Jennings, C., Bostr\"{o}m, H., and J-I. Bruaroey, ``WebRTC 1.0:
  Real-time Communication Between Browsers'', W3C Proposed
  Recommendation, \url{https://www.w3.org/TR/webrtc/}.

\end{description}

\subsection{Informative References}

\begin{description}[style=nextline, leftmargin=2cm, labelwidth=1.8cm]

\item[{[fetch]}] van Kesteren, A., ``Fetch'',
  \url{https://fetch.spec.whatwg.org/}.

\item[{[RFC3261]}] Rosenberg, J., Schulzrinne, H., Camarillo, G., Johnston,
  A., Peterson, J., Sparks, R., Handley, M., and E.
  Schooler, ``SIP: Session Initiation Protocol'', RFC 3261,
  DOI 10.17487/RFC3261, June 2002,
  \url{https://www.rfc-editor.org/info/rfc3261}.

\item[{[RFC5705]}] Rescorla, E., ``Keying Material Exporters for Transport
  Layer Security (TLS)'', RFC 5705, DOI 10.17487/RFC5705,
  March 2010, \url{https://www.rfc-editor.org/info/rfc5705}.

\item[{[RFC6120]}] Saint-Andre, P., ``Extensible Messaging and Presence
  Protocol (XMPP): Core'', RFC 6120, DOI 10.17487/RFC6120,
  March 2011, \url{https://www.rfc-editor.org/info/rfc6120}.

\item[{[RFC6265]}] Barth, A., ``HTTP State Management Mechanism'', RFC 6265,
  DOI 10.17487/RFC6265, April 2011,
  \url{https://www.rfc-editor.org/info/rfc6265}.

\item[{[RFC6455]}] Fette, I. and A. Melnikov, ``The WebSocket Protocol'',
  RFC 6455, DOI 10.17487/RFC6455, December 2011,
  \url{https://www.rfc-editor.org/info/rfc6455}.

\item[{[RFC6943]}] Thaler, D., Ed., ``Issues in Identifier Comparison for
  Security Purposes'', RFC 6943, DOI 10.17487/RFC6943, May
  2013, \url{https://www.rfc-editor.org/info/rfc6943}.

\item[{[RFC7617]}] Reschke, J., ``The `Basic' HTTP Authentication Scheme'',
  RFC 7617, DOI 10.17487/RFC7617, September 2015,
  \url{https://www.rfc-editor.org/info/rfc7617}.

\item[{[RFC8224]}] Peterson, J., Jennings, C., Rescorla, E., and C. Wendt,
  ``Authenticated Identity Management in the Session
  Initiation Protocol (SIP)'', RFC 8224,
  DOI 10.17487/RFC8224, February 2018,
  \url{https://www.rfc-editor.org/info/rfc8224}.

\item[{[RFC8828]}] Uberti, J. and G. Shieh, ``WebRTC IP Address Handling
  Requirements'', RFC 8828, DOI 10.17487/RFC8828, January
  2021, \url{https://www.rfc-editor.org/info/rfc8828}.

\item[{[TLS-DTLS13]}] Rescorla, E., Tschofenig, H., and N. Modadugu, ``The
  Datagram Transport Layer Security (DTLS) Protocol Version
  1.3'', Work in Progress, Internet-Draft, draft-ietf-tls-dtls13-39, 2 November 2020,
  \url{https://tools.ietf.org/html/draft-ietf-tls-dtls13-39}.

\end{description}

\section*{Acknowledgements}

Bernard Aboba, Harald Alvestrand, Richard Barnes, Dan Druta, Cullen
Jennings, Hadriel Kaplan, Matthew Kaufman, Jim McEachern, Martin
Thomson, Magnus Westerlund.  Matthew Kaufman provided the UI material
in Section 6.5.  Christer Holmberg provided the initial version of
Section 5.1.

\section*{Author's Address}

Eric Rescorla\\
Mozilla

Email: \href{mailto:ekr@rtfm.com}{ekr@rtfm.com}

\end{document}
